\section{Probabilistic Graphical Models: Belief Networks and Parameter Learning}
\label{sec:pgm-fundamentals}

\subsection{Belief Networks}

A belief network (also called a Bayesian network) represents the joint
probability distribution of a set of random variables $X_1, \ldots, X_m$ as a
directed acyclic graph (DAG) $G = \langle V, E \rangle$, where each node
$v \in V$ corresponds to one variable $X_v$, and each directed edge encodes a
direct probabilistic dependency. The defining property of this representation
is that the joint density factorises along the graph structure:
\[
  p(x_1, \ldots, x_m) = \prod_{v=1}^{m} p\bigl(x_v \mid \mathrm{pa}(x_v)\bigr),
\]
where $\mathrm{pa}(x_v)$ denotes the values taken by the parents of $v$ in
$G$ (and $p(x_v \mid \mathrm{pa}(x_v)) = p(x_v)$ for a node with no parents).
This factorisation is what makes belief networks tractable: instead of
specifying one joint distribution over all $m$ variables jointly, the modeller
only needs to specify $m$ local conditional distributions, one per
node, each depending only on that node's (typically small) parent set
\autocite{pearl1988probabilistic,koller2009pgm}.

A standard illustrative example is the "Rain--Sprinkler--Wet Grass" network:
whether the grass is wet depends on whether it rained and whether the
sprinkler was on, and both of those, in turn, have their own (unconditional)
probabilities of occurring. Once the network and its conditional probability
tables are specified, exact inference lets us answer queries in the reverse
causal direction — e.g.\ given that we observe the grass is wet, what
is the probability that it rained? — via Bayes' theorem and marginalisation
over the remaining variables. The same machinery applies, at much larger
scale, to our 24-variable oil-price panel: observing a change in, say, OPEC
spare capacity and the US dollar index lets the network update its belief
about the probability distribution of the oil-price node, without needing a
separate model for every possible combination of observed evidence.

\subsection{Learning a Belief Network}

Given a dataset of observations but no pre-specified network, learning a
belief network splits into two distinct sub-problems:

\begin{description}
  \item[Parameter learning] Given a fixed DAG $G$ that already captures the
    dependency structure between variables, estimate the conditional
    probability distribution (CPD) at each node, i.e.\ estimate the
    parameters $\theta_{v \mid \mathrm{pa}(v)}$ of
    $p_{v \mid \mathrm{pa}(v)}(x_v; \theta_{v\mid\mathrm{pa}(v)})$.
  \item[Structure learning] Given only the raw data, estimate the DAG $G$
    itself, i.e.\ decide which arcs are present, without yet committing to
    any parameter values. This is covered in
    Section~\ref{sec:structure-learning}.
\end{description}

This section focuses on parameter learning. One of the principal practical
advantages of the graphical representation is that it decomposes one
high-dimensional estimation problem (the full joint distribution over all
$m$ variables) into $m$ independent, low-dimensional estimation problems, one
per node, each conditioned only on that node's parents. This is only
computationally cheap if the graph is reasonably sparse: the number of free
parameters at a node grows with the number (and cardinality) of its parents,
so there is a direct computational incentive to prefer sparser structures —
a consideration that resurfaces in Section~\ref{sec:structure-learning} when
scoring candidate structures.

Two standard estimators are used for the CPD parameters at each node:

\begin{description}
  \item[Maximum Likelihood Estimation (MLE).] The CPD at each node is filled in
    using the observed relative frequencies of the variable's states (for
    discrete data), choosing the parameter values that maximise
    $P(\text{data} \mid \text{model})$. This is simple and consistent, but can
    behave poorly when some parent-value combinations are rarely or never
    observed in the sample (a combination with zero observed occurrences
    receives zero probability, which is often too strong a conclusion to draw
    from a finite sample).
  \item[Bayesian Estimation.] Rather than starting from an uninformative
    (empirical-frequency) prior, Bayesian parameter estimation begins from an
    explicit prior CPD expressing the modeller's belief about the parameters
    before observing the data, and updates it with the observed data
    via Bayes' rule. This naturally regularises the rare/unseen
    parent-combination problem that pure MLE struggles with, at the cost of
    requiring the modeller to choose a (possibly subjective) prior.
\end{description}

\subsection{Relevance to this Project}

In our oil-price panel, the "variables" that ultimately enter the belief network are not the raw, continuous series in Table~\ref{tab:data-dictionary}, but discretised regime labels produced by the Hidden Markov Model stage (e.g.\
"WTI: bull/stagnant/bear", "USD index: strengthening/stable/weakening"). Once each series has been reduced to a small number of discrete regimes, parameter
learning reduces to exactly the discrete-CPD estimation problem described above, and can be carried out directly with standard tools (e.g.\ \texttt{pgmpy}'s
maximum-likelihood or Bayesian estimators \autocite{ankan2015pgmpy}) once the
network structure itself has been learned, as discussed next.
